% Generated by tools/edn_latex.py from reports/edn.md.
% Do not edit: change reports/edn.md and run `make edn`.
\ednentry{
  id = {EDN-75},
  title = {The test suite also runs on Windows in CI, as a required check},
  date = {2026-10-06},
  milestone = {M3: Quality Assurance},
  topic = {Testing Strategy, CI/CD},
  participants = {@lukas2510, decided by the repository owner; the Windows failures that prompted it were found by @kadameit},
  decision = {CI runs the whole test suite on \texttt{windows-\allowbreak{}latest} as well as on Linux, for every change that can affect it, as the job \texttt{Tests (Windows)}, and the job is a required check like \texttt{Tests}.
Coverage and the test reports still come from the Linux job alone.},
  rationale = {\emph{Alternatives considered:}
\begin{itemize}[leftmargin=*, topsep=2pt]
\item \textbf{Option A (chosen): a required Windows job.}
\newline Pros: a Windows-only failure is found on Windows, before the merge, by whoever opens the pull request, instead of after it by whoever happens to develop on Windows; it covers every failure mode the suite can reach, including ones nobody has thought of yet; and it costs nothing in money, because the repository is public.
\newline Cons: every code pull request waits for a second runner, which took 374 s against the Linux job's 214 s on the pull request that added it; a test that is flaky only on Windows blocks every merge until it is fixed; and a test that starts processes or touches paths has to be written for both platforms, as the local MLflow server of \texttt{tests/\allowbreak{}test\_\allowbreak{}tracking.\allowbreak{}py} had to be stopped with \texttt{taskkill /\allowbreak{}T} on Windows and by process group on POSIX.
\item \textbf{Option B: the same job, not required.}
\newline Pros: a Windows-only flake can never block a merge.
\newline Cons: a red check that does not block is a red check people learn to merge past, so a regression still reaches \texttt{main} and is found where it is found today, by the next person on Windows.
\item \textbf{Option C: emulate Windows in the Linux suite.}
Tests that run with a cp1252 stdout, or that patch text-mode writes to translate to CRLF.
\newline Pros: no second runner and no added wait.
\newline Cons: an emulation covers only the failure modes it was written for.
It would not have found the cross-drive failure below, which nobody knew of, nor \#74, and the CRLF part could only be emulated by patching \texttt{Path.\allowbreak{}write\_\allowbreak{}text}, which tests the patch rather than the platform.
\item \textbf{Option D: keep relying on contributors who develop on Windows to run the suite there.}
\newline Pros: no cost at all.
\newline Cons: it is how \#73, \#74 and \#84 were found, each after the code had merged, and it depends on one person running the suite and reading the result.
\end{itemize}
\par
\emph{Why:} one of us develops on Windows, so the project has to work there, and NFR-06 asks that a clean clone reproduces the pipeline, which on Windows it did not: every JSON artefact came out with CRLF line endings and a different hash, and the MLflow setup check of the getting-started page crashed on an emoji it could not print (\#84).
Only a job on the platform itself catches such failures without knowing them in advance, and only a required one keeps them off \texttt{main}.
\par
The first run of the job made the case on its own.
On the commit that added the two tests for \#84 before the fix, Linux failed one test, the setup check against a local MLflow server with a cp1252 stdout, while Windows failed seventeen: that test, the line-ending test, which cannot fail on Linux, and fifteen tests that failed with \texttt{ValueError: path is on mount \textquotesingle{}D:\textquotesingle{}, start on mount \textquotesingle{}C:\textquotesingle{}}.
GitHub's Windows runners check the repository out on \texttt{D:} and keep the temporary directory on \texttt{C:}, and DVC resolves \texttt{dvc.\allowbreak{}yaml} relative to the current directory, which has no relative form across drives.
The same failure hits any contributor whose clone and current directory are on different drives, and it turned tracking off in \texttt{train}: \texttt{recommenditos/\allowbreak{}provenance.\allowbreak{}py} now reads a DVC repository with its root as the current directory.
After the fixes, both jobs pass the same 679 tests.
\par
The wait was accepted: the Windows job runs in parallel with the others, so a code pull request waits about two and a half minutes longer for its last required check, 374 s instead of the Linux job's 214 s.},
  ai = {Information seeking, Alternative generation, Alternative assessment, Recommendation, Solution generation},
  response = {Accepted},
  assessment = {AI reviewed and merged Kevin's \#73 and \#75, then traced the two Windows failures \#73 left open to their causes: the seven JSON writes that relied on the platform's line separator, and MLflow's own \texttt{sys.\allowbreak{}stdout.\allowbreak{}write} of the run link after an emoji, which also fails a stage under \texttt{RECOMMENDITOS\_\allowbreak{}REQUIRE\_\allowbreak{}TRACKING=\allowbreak{}1}.
It reproduced the console failure end to end before fixing it, by running the setup check against DagsHub with \texttt{PYTHONIOENCODING=\allowbreak{}cp1252}, generated options A to C with their pros and cons, recommended A, and Lukas chose A as recommended; option D is the status quo, added here for completeness.
The fixes and the tests, written to fail first, were AI's, and so was the analysis of the fifteen cross-drive failures the new job found.},
  aievidence = {Claude Code session of 2026-10-06: the review of \#73 and \#75, in which the follow-ups were found, and the work on issue \#84, in which the setup check was reproduced against DagsHub with a cp1252 stdout (exit 1, \texttt{UnicodeEncodeError: \textquotesingle{}charmap\textquotesingle{} codec can\textquotesingle{}t encode character \textquotesingle{}\textbackslash{}U0001f3c3\textquotesingle{}}) and the options were put to Lukas.},
  evidence = {\href{https://github.com/mlops-2627q1-mds-upc/MLOPS_Recommenditos/issues/84}{issue \#84}; \href{https://github.com/mlops-2627q1-mds-upc/MLOPS_Recommenditos/pull/85}{PR \#85}; the \href{https://github.com/mlops-2627q1-mds-upc/MLOPS_Recommenditos/actions/runs/37438994409}{failing run} on the tests before the fix and the \href{https://github.com/mlops-2627q1-mds-upc/MLOPS_Recommenditos/actions/runs/37440411226}{passing run} after it; \href{https://github.com/mlops-2627q1-mds-upc/MLOPS_Recommenditos/pull/73}{\#73}, \href{https://github.com/mlops-2627q1-mds-upc/MLOPS_Recommenditos/issues/74}{\#74}; \texttt{.\allowbreak{}github/\allowbreak{}workflows/\allowbreak{}ci.\allowbreak{}yml}, job \texttt{test-\allowbreak{}windows}; \href{https://github.com/mlops-2627q1-mds-upc/MLOPS_Recommenditos/blob/main/CONTRIBUTING.md}{Contributing}, Before opening a PR.},
}
